% F082｜02 理解
\clearpage
\begin{explanationbox}
	\textbf{\textcolor{CExplanation}{一、为什么会有“距离差”这种解法？}}
\begin{center}\begin{tikzpicture}[font=\small,>=stealth,line join=round]
\draw[gray!40,rounded corners=3pt] (0,0) rectangle (14.2,10.6);\node at (7.1,10.0) {两个根号相减：蓝线是橙、紫两条曲线的高度差};
\node[blue] at (7.1,9.1) {$f(x)=\sqrt{x^2+4}-\sqrt{(x-2)^2+1}$};\begin{scope}[shift={(3.0,2.9)},x=.85cm,y=.65cm]\draw[->,gray] (-2.2,0)--(6,0) node[right]{$x$};\draw[->,gray] (0,-1.7)--(0,7.7) node[left]{$y$};\foreach \v in {-1,2,4,6}{\draw[gray] (-.06,\v)--(.06,\v) node[left=3pt]{$\v$};}\foreach \v in {2,3,4}{\draw[gray] (\v,-.07)--(\v,.07) node[below=4pt]{$\v$};}\draw[orange!85!black,thick,dashed,domain=-2:5,samples=150] plot(\x,{sqrt(\x*\x+4)});\draw[violet,thick,densely dotted,domain=-2:5,samples=150] plot(\x,{sqrt((\x-2)^2+1)});\draw[blue,very thick,domain=-2:5,samples=180] plot(\x,{sqrt(\x*\x+4)-sqrt((\x-2)^2+1)});\draw[gray,dotted] (3,0)--(3,{sqrt(13)});\fill[orange!85!black] (3,{sqrt(13)}) circle(2pt);\fill[violet] (3,{sqrt(2)}) circle(2pt);\fill[blue] (3,{sqrt(13)-sqrt(2)}) circle(2pt);\draw[gray,dotted] (3,{sqrt(2)})--(5.6,{sqrt(2)}) (3,{sqrt(13)})--(5.6,{sqrt(13)});\draw[blue,<->] (5.6,{sqrt(2)})--(5.6,{sqrt(13)}) node[midway,right=4pt]{$f(3)$};\end{scope}\node[anchor=west,align=left] at (9.3,6.3) {\textcolor{orange!85!black}{$U=\sqrt{x^2+4}=PA$}\\[13pt]\textcolor{violet}{$V=\sqrt{(x-2)^2+1}=PB$}\\[13pt]\textcolor{blue}{$f=U-V$：高度差}\\[13pt]第一项更小，差就为负。};
\node[text width=12.8cm,align=center] at (7.1,.65) {在同一个 $x=3$ 处相减；蓝线高度为 $\sqrt{13}-\sqrt2$，不是两项的最低值相减。};
\end{tikzpicture}\end{center}



	先看一个根号：
	\[
		\sqrt{(x-m)^2+a^2}.
	\]
	把 $P(x,0)$ 放在水平的 $x$ 轴上，把 $A(m,a)$ 放在轴上方。水平方向相差 $|x-m|$，竖直方向相差 $a$。按勾股定理，两点间的直线距离就是这个根号。

	另一个根号也是如此：取 $B(n,b)$，于是
	\[
		\underbrace{\sqrt{(x-m)^2+a^2}}_{PA}
		-\underbrace{\sqrt{(x-n)^2+b^2}}_{PB}
		=PA-PB.
	\]
	\textbf{加法}表示两段距离的\textbf{总长度}；\textbf{减法}表示\textbf{第一段比第二段长多少}，可能是正数、负数或零。这里不必把差理解成真实走路的总长度。

	看蓝线落在横轴上方还是下方，就能看到距离差为正还是为负。\textbf{方向由相减的先后顺序决定。}

	\medskip
\end{explanationbox}
\clearpage
\begin{explanationbox}
	\textbf{\textcolor{CExplanation}{二、为什么差的绝对值不能超过 $AB$？}}
\begin{center}\begin{tikzpicture}[font=\small,>=stealth,line join=round]
\draw[gray!40,rounded corners=3pt] (0,0) rectangle (14.2,8.4);\node at (7.1,7.800000000000001) {比较的是直路 PA 与绕路 PBA，先得到范围界};
\begin{scope}[shift={(3.0,2.2)},x=1.2cm,y=1.2cm]\draw[->,gray] (-1.6,0)--(2.6,0) node[right]{$x$};\draw[->,gray] (0,-.35)--(0,2.8) node[left]{$y$};\draw[orange!85!black,line width=1.6pt] (0,2)--(-1,0);\draw[violet,thick] (2,1)--(-1,0);\draw[teal!80!black,thick] (0,2)--(2,1);\fill[teal!80!black] (0,2) circle(2.2pt) node[above left=4pt]{$A(0,2)$};\fill[teal!80!black] (2,1) circle(2.2pt) node[above right=4pt]{$B(2,1)$};\fill[blue] (-1,0) circle(2.2pt) node[below=4pt]{$P(-1,0)$};\end{scope}\node[anchor=west,align=left] at (8.6,4.4) {直路：$PA$\\[12pt]绕行：$PB+BA$\\[12pt]$PA\le PB+BA$\\[12pt]$PA-PB\le AB$\\[12pt]交换 $A,B$ 再得另一侧界。};
\node[text width=12.8cm,align=center] at (7.1,.65) {橙色为直路；紫色与青色合起来是绕路。这里只得到界，还未证明能取到。};
\end{tikzpicture}\end{center}



	看图比较从 $P$ 直接到 $A$ 与先到 $B$ 再到 $A$：直路 $PA$ 不会长于绕路 $PB+BA$，因此
	\[
		PA\leq AB+PB\quad\Longrightarrow\quad PA-PB\leq AB.
	\]
	交换 $A,B$ 的位置看，也有
	\[
		PB\leq AB+PA\quad\Longrightarrow\quad PB-PA\leq AB.
	\]
	两句话合起来，就是
	\[
		\boxed{-AB\leq PA-PB\leq AB}.
	\]
	这说的是：\textbf{站在同一个 $P$，到两处地点的路程差，不可能比两处地点本身的距离还大。}

	但要记住：\textbf{“不超过”不等于“一定能等于”。} 下一步必须看图找等号。
\end{explanationbox}
\clearpage
\begin{explanationbox}
	\textbf{\textcolor{CExplanation}{三、等号什么时候才能成立？}}

	若 $PA-PB=AB$，则 $PA=PB+BA$。这表示从 $P$ 走到 $B$，再走到 $A$，刚好是一条不折返的直路，即\textbf{$P,B,A$ 依次排成一条直线}。

	若 $PA-PB=-AB$，同理必须是\textbf{$P,A,B$ 依次排成一条直线}。

	因此求\textbf{距离差}，并不是让 $P$ 处在 $A,B$ 中间；相反，\textbf{$P$ 必须在两点连线的外侧}。这和 F012 距离和型“$P$ 在线段中间时最短”正好相反。

	下面的图取 $A(0,2)$、$B(2,1)$、$P(4,0)$，三点恰好在一条直线上，且 $P$ 在 $B$ 的外侧。
	\begin{center}\begin{tikzpicture}[font=\small,>=stealth,line join=round]
\draw[gray!40,rounded corners=3pt] (0,0) rectangle (14.2,8.6);\node at (7.1,8.0) {延长线上的等号：长距离减短距离，剩下两定点的距离};
\begin{scope}[shift={(1.9,2.2)},x=1.2cm,y=1.2cm]\draw[->,gray] (-0.6,0)--(4.6,0) node[right]{$x$};\draw[->,gray] (0,-.35)--(0,2.8) node[left]{$y$};\draw[orange!85!black,line width=1.6pt] (0,2)--(4,0);\draw[violet,thick] (2,1)--(4,0);\draw[teal!80!black,thick] (0,2)--(2,1);\fill[teal!80!black] (0,2) circle(2.2pt) node[above left=4pt]{$A(0,2)$};\fill[teal!80!black] (2,1) circle(2.2pt) node[above right=4pt]{$B(2,1)$};\fill[blue] (4,0) circle(2.2pt) node[below=4pt]{$P(4,0)$};\end{scope}\node[anchor=west,align=left] at (8.4,4.5) {$A(0,2),\ B(2,1)$\\[11pt]$P(4,0)$ 在 $AB$ 外侧\\[11pt]\textcolor{orange!85!black}{$PA=2\sqrt5$}\\[10pt]\textcolor{violet}{$PB=\sqrt5$}\\[10pt]\textcolor{blue}{$PA-PB=AB=\sqrt5$}};
\node[text width=12.8cm,align=center] at (7.1,.65) {图中横纵同尺度；三点顺序为 $A,B,P$，取到的是正端点。};
\end{tikzpicture}\end{center}

	\[
		PA=2\sqrt5,\quad PB=\sqrt5,\quad AB=\sqrt5,
		\qquad PA-PB=AB.
	\]
	所以 $\sqrt5$ 不只是一个算出来的上界，\textbf{它真的能取到}。
\end{explanationbox}
\clearpage
\begin{explanationbox}
	\textbf{\textcolor{CExplanation}{四、最值得深入：为什么“翻折一个定点”有用？}}

	根号里竖直距离是平方：
	\[
		(0-a)^2=(0+a)^2.
	\]
	因此，$A(m,a)$ 与 $A'(m,-a)$ 到任意 $P(x,0)$ 的距离完全一样。\textbf{把一个固定点关于 $x$ 轴翻折，不改变根号的值。}

	\begin{center}\begin{tikzpicture}[font=\small,>=stealth,line join=round]
\draw[gray!40,rounded corners=3pt] (0,0) rectangle (14.2,9.2);\node at (7.1,8.6) {翻折不改变某一根号，却会改变两定点之间的距离};
\begin{scope}[shift={(1.9,4.1)},x=1.0cm,y=1.0cm]\draw[->,gray] (-0.6,0)--(4.6,0) node[right]{$x$};\draw[->,gray] (0,-.35)--(0,2.8) node[left]{$y$};\draw[orange!85!black,line width=1.6pt] (0,2)--(4,0);\draw[violet,thick] (2,1)--(4,0);\draw[teal!80!black,thick] (0,2)--(2,1);\fill[teal!80!black] (0,2) circle(2.2pt) node[above left=4pt]{$A(0,2)$};\fill[teal!80!black] (2,1) circle(2.2pt) node[above right=4pt]{$B(2,1)$};\fill[blue] (4,0) circle(2.2pt) node[below=4pt]{$P(4,0)$};\end{scope}\begin{scope}[shift={(1.9,4.1)}]\draw[gray,dashed] (0,-2)--(4,0);\draw[gray,dotted] (0,-2)--(0,2);\draw[gray,thick,dashed] (0,-2)--(2,1);\fill[gray] (0,-2) circle(2pt) node[below right=3pt]{$A^{\prime}(0,-2)$};\end{scope}\node[anchor=west,align=left] at (8.2,4.5) {$PA^{\prime}=PA$\\[12pt]但 $AB=\sqrt5$\\[10pt]而 $A^{\prime}B=\sqrt{13}$\\[12pt]求差用同侧点的 $AB$。};
\node[text width=12.8cm,align=center] at (7.1,.65) {橙、灰两段到轴上同一点的距离相等；同侧和异侧的定点距离不能混用。};
\end{tikzpicture}\end{center}


	从上面可以看到：同样的两个根号，可以画成上下两侧，也可以翻折成同一侧。\textbf{对于距离差，我们选择两个固定点在 $x$ 轴同一侧}，因为要找的是$P$落在两点连线的\textbf{延长线}上时，距离差能否等于 $AB$。

	注意这个容易混淆的地方：\textbf{翻折后，$PA$ 不变，但 $AB$ 可能改变。} 所以不能在点的上下位置随便改变之后，把\textbf{旧图的 $AB$} 继续当作新的边界；必须按最终选定的两个同侧点，重新计算它们之间的距离：
	\[
		AB=\sqrt{(m-n)^2+(a-b)^2}.
	\]
	两个定点放\textbf{同侧}，求差用的是高度差 $|a-b|$；F012 中两点放\textbf{异侧}、求和用的是高度和 $a+b$。
\end{explanationbox}
\clearpage
\begin{explanationbox}
	\textbf{\textcolor{CExplanation}{五、高度不同，为什么能找到等号位置？}}
\begin{center}\begin{tikzpicture}[font=\small,>=stealth,line join=round]
\draw[gray!40,rounded corners=3pt] (0,0) rectangle (14.2,8.0);\node at (7.1,7.4) {高低换过来，延长线的交点与距离差的正负也改变};
\begin{scope}[shift={(2.0,3.2)},x=0.75cm,y=0.75cm]\draw[->,gray] (-0.6,0)--(4.6,0) node[right]{$x$};\draw[->,gray] (0,-.35)--(0,2.8) node[left]{$y$};\draw[orange!85!black,line width=1.6pt] (0,2)--(4,0);\draw[violet,thick] (2,1)--(4,0);\draw[teal!80!black,thick] (0,2)--(2,1);\fill[teal!80!black] (0,2) circle(2.2pt) node[left=6pt]{$A$};\fill[teal!80!black] (2,1) circle(2.2pt) node[above right=4pt]{$B$};\fill[blue] (4,0) circle(2.2pt) node[below=4pt]{$P$};\end{scope}\node at (3.2,6.3) {$a>b$：正端点};\node at (3.2,1.65) {$f(4)=\sqrt5$};\begin{scope}[shift={(10.0,3.2)},x=0.75cm,y=0.75cm]\draw[->,gray] (-2.6,0)--(2.6,0) node[right]{$x$};\draw[->,gray] (0,-.35)--(0,2.8) node[left]{$y$};\draw[orange!85!black,line width=1.6pt] (0,1)--(-2,0);\draw[violet,thick] (2,2)--(-2,0);\draw[teal!80!black,thick] (0,1)--(2,2);\fill[teal!80!black] (0,1) circle(2.2pt) node[left=6pt]{$A$};\fill[teal!80!black] (2,2) circle(2.2pt) node[above right=4pt]{$B$};\fill[blue] (-2,0) circle(2.2pt) node[below=4pt]{$P$};\end{scope}\node at (11.2,6.3) {$a<b$：负端点};\node at (11.2,1.65) {$f(-2)=-\sqrt5$};
\node[text width=12.8cm,align=center] at (7.1,.65) {同尺度对照：$a\ne b$ 时 $x_0=(bm-an)/(b-a)$；交点都在线段外侧。};
\end{tikzpicture}\end{center}



	先令 $a>b$：$A(m,a)$ 比 $B(n,b)$ 高一些。连接 $A,B$ 并继续朝着较低的 $B$ 方向延长，高度从 $a$ 下降到 $b$，最后一定降到 $0$，与 $x$ 轴相交。交点就是 $P_0$。

	此时顺序为 $A,B,P_0$，所以 $P_0A=P_0B+BA$，原式最大值为 $AB$。

	当 $a<b$ 时，换成沿着 $A$ 的方向延长，交点落在 $A$ 一侧，得到 $P_0B=P_0A+AB$，故原式最小值为 $-AB$。

	\textbf{怎么求交点？} 直线从 $A(m,a)$ 走到 $B(n,b)$ 时，横坐标改变量为 $n-m$，高度改变量为 $b-a$。继续沿同一直线走到高度 $0$，因此
	\[
		\frac{x_0-m}{n-m}=\frac{-a}{b-a}\qquad(n\ne m),
	\]
	整理得到\textbf{不必再分 $m,n$ 大小}的统一公式：
	\[
		\boxed{x_0=\frac{bm-an}{b-a}\quad(a\ne b)}.
	\]
	若 $m=n$，两点在同一条竖直线上，延长线与轴交于 $x_0=m$；这个公式同样成立。

	\textbf{一个很实用的判断：} $a>b$ 看最大值 $D$；$a<b$ 看最小值 $-D$。另一个边界需要另外检查，绝不能顺手写成 $[-D,D]$。
\end{explanationbox}
\clearpage
\begin{explanationbox}
	\textbf{\textcolor{CExplanation}{六、高度相同，为什么反而取不到两侧边界？}}

	若 $a=b>0$，则 $A(m,a)$、$B(n,a)$ 一样高。两点连线是一条\textbf{水平直线}，与 $x$ 轴平行。
	\begin{center}\begin{tikzpicture}[font=\small,>=stealth,line join=round]
\draw[gray!40,rounded corners=3pt] (0,0) rectangle (14.2,8.0);\node at (7.1,7.4) {两定点一样高：延长线与 x 轴永远不相交};
\begin{scope}[shift={(2.7,2.3)},x=1cm,y=1cm]\draw[->,gray] (-1.6,0)--(3.6,0) node[right]{$x$};\draw[->,gray] (0,-.35)--(0,2.8) node[left]{$y$};\draw[orange!85!black,line width=1.6pt] (0,2)--(-1,0);\draw[violet,thick] (3,2)--(-1,0);\draw[teal!80!black,thick] (0,2)--(3,2);\fill[teal!80!black] (0,2) circle(2.2pt) node[above left=4pt]{$A(0,2)$};\fill[teal!80!black] (3,2) circle(2.2pt) node[above right=4pt]{$B(3,2)$};\fill[blue] (-1,0) circle(2.2pt) node[below=4pt]{$P(-1,0)$};\end{scope}\draw[teal!80!black,dashed] (1.2,4.3)--(8,4.3);\node[anchor=west,align=left] at (8.6,4.4) {$A(0,2),\ B(3,2)$\\[12pt]$AB$ 所在直线为 $y=2$\\[12pt]与动点所在 $y=0$ 平行\\[12pt]不存在取等所需的交点。};
\node[text width=12.8cm,align=center] at (7.1,.65) {只要两定点不同且都在轴上方，距离差的两端都不能取到。};
\end{tikzpicture}\end{center}


	因为 $P$ 必须在 $x$ 轴上，不论移动到哪里，都不可能与 $A,B$ 排成所需的同一直线（除非两点恰好重合，此时另算）。因此，当 $m\ne n$ 时：
	\[
		\boxed{-|m-n|<f(x)<|m-n|}.
	\]
	把 $P$ 一直往右移，两段斜边都会接近它们各自的水平长度，两者的差会越来越接近 $n-m$；一直往左移，就越来越接近 $m-n$。\textbf{靠近边界，不表示取到边界。}

	例如
	\[
		f(x)=\sqrt{x^2+4}-\sqrt{(x-3)^2+4}
	\]
	的两个固定点是 $(0,2)$、$(3,2)$，所以值域是 $(-3,3)$，而\textbf{不是} $[-3,3]$。

	若再有 $m=n$，两个根号从头到尾都一样，$f(x)\equiv0$；此时值域只是 $\{0\}$。
\end{explanationbox}
\clearpage
\begin{explanationbox}
	\textbf{\textcolor{CExplanation}{七、另一个边界从哪里来？}}
\begin{center}\begin{tikzpicture}[font=\small,>=stealth,line join=round]
\draw[gray!40,rounded corners=3pt] (0,0) rectangle (14.2,10.2);\node at (7.1,9.6) {第一项高度更大：先增后减，只有最大值能取到};
\node[blue] at (7.1,8.65) {$\displaystyle f(x)=\sqrt{x^2+4}-\sqrt{(x-2)^2+1}$};\begin{scope}[shift={(3.1,4.6)},x=0.36cm,y=1cm]\draw[->,gray] (-6,0)--(14.5,0) node[right]{$x$};\draw[->,gray] (0,-2.8)--(0,3.4000000000000004) node[left]{$y$};\draw[gray,dashed] (-6,-2)--(14,-2);\node[left=3pt,fill=white,inner sep=1pt] at (0,-2) {$-2$};\draw[gray,dashed] (-6,2)--(14,2);\node[left=3pt,fill=white,inner sep=1pt] at (0,2) {$2$};\draw[blue,very thick,domain=-6:14,samples=240] plot(\x,{sqrt(\x*\x+4)-sqrt((\x-2)^2+1)});\draw[gray,dotted] (4,0)--(4,2.23606797749979);\fill[blue] (4,2.23606797749979) circle(2.3pt);\node[below=4pt] at (4,0){$4$};\node[blue,above=9pt] at (4,2.23606797749979) {$(4,\sqrt5)$};\end{scope}\node[anchor=west,align=left] at (9,4.5) {$a=2>b=1$\\[12pt]左端：趋近 $-2$\\[10pt]右端：趋近 $2$\\[12pt]最大值：$\sqrt5$\\[12pt]完整值域：$(-2,\sqrt5]$};
\node[text width=12.8cm,align=center] at (7.1,.65) {虚线是无限远处靠近的高度；不是有限定义域的端点。};
\end{tikzpicture}\end{center}



	当 $x$ 很大且为正数时，两个根号分别非常接近 $x-m$、$x-n$，所以差接近 $n-m$；当 $x$ 是绝对值很大的负数时，两个根号分别接近 $m-x$、$n-x$，所以差接近 $m-n$：
	\[
		\boxed{\begin{aligned}
				x\to+\infty & :\quad f(x)\to n-m, \\
				x\to-\infty & :\quad f(x)\to m-n.
			\end{aligned}}
	\]
	这两个数中较小的是 $-d$，较大的是 $d$。因此，$a>b$ 时，函数在一侧达到真正最大值 $D$，另一端无限接近 $-d$；$a<b$ 时则反过来。

	\textbf{为什么不能把 $-d$ 或 $d$ 写成闭端点？} 当 $a>b$ 时，第一个固定点更高，第一段距离总比把它压到高度 $b$ 后长，故函数值\textbf{严格大于} $-d$；$a<b$ 时同理严格小于 $d$。完整计算在“证明二”。

	\textbf{对高考题最快的提醒：}先分 $a>b$、$a<b$、$a=b$；再看边界是\textbf{某个有限的 $x$ 真正取到}，还是\textbf{$x$ 无限远时才慢慢靠近}。取到写中括号；只靠近写小括号。
\end{explanationbox}
\clearpage
\begin{explanationbox}
	\textbf{\textcolor{CExplanation}{八、再发现一个简单却有用的关系：差为零在哪里？}}
\begin{center}\begin{tikzpicture}[font=\small,>=stealth,line join=round]
\draw[gray!40,rounded corners=3pt] (0,0) rectangle (14.2,8.7);\node at (7.1,8.1) {例8：根号差为零，就是到两定点一样远};
\begin{scope}[shift={(2.5,2.2)},x=1.5cm,y=1.5cm]\draw[->,gray] (-0.6,0)--(2.6,0) node[right]{$x$};\draw[->,gray] (0,-.35)--(0,2.8) node[left]{$y$};\draw[orange!85!black,line width=1.6pt] (0,2)--(0.25,0);\draw[violet,thick] (2,1)--(0.25,0);\draw[teal!80!black,thick] (0,2)--(2,1);\fill[teal!80!black] (0,2) circle(2.2pt) node[above left=4pt]{$A$};\fill[teal!80!black] (2,1) circle(2.2pt) node[above right=4pt]{$B$};\fill[blue] (0.25,0) circle(2.2pt) node[below=4pt]{$P$};\end{scope}\begin{scope}[shift={(2.5,2.2)},x=1.5cm,y=1.5cm]\draw[gray,dashed] (.1,-.3)--(1.55,2.6);\fill[gray] (1,1.5) circle(1.8pt);\node[above left=7pt] at (1,1.5){$M$};\draw[gray] (1.134,1.433)--(1.201,1.567)--(1.067,1.634);\end{scope}\node[anchor=west,align=left] at (8.5,4.4) {$PA=PB$\\[12pt]$P$ 在 $AB$ 的垂直平分线上\\[12pt]与 $x$ 轴交于 $x=1/4$\\[12pt]有理化后分子为 $4x-1$\\[10pt]所以零点仅为 $1/4$。};
\node[text width=12.8cm,align=center] at (7.1,.65) {橙、紫两段等长；灰点 $M$ 为 $AB$ 中点，灰虚线为垂直平分线。};
\end{tikzpicture}\end{center}



	两段距离恰好相等，即 $f(x)=0$ 时，有 $PA=PB$。两个正根号相等可以直接平方：
	\[
		(x-m)^2+a^2=(x-n)^2+b^2.
	\]
	平方项展开，$x^2$ 正好抵消，剩下一次式：
	\[
		2(n-m)x+(m^2-n^2)+(a^2-b^2)=0.
	\]
	当 $m\ne n$ 时，交点横坐标为
	\[
		\boxed{x=\frac{m+n}{2}+\frac{b^2-a^2}{2(n-m)}}.
	\]
	若 $m=n$ 且 $a\ne b$，两段距离始终不相等；若 $m=n$ 且 $a=b$，则处处相等。

	\textbf{眼前一亮：}距离差看上去有两个根号，\textbf{判断正负或求零点时，根号却可以一起消掉}。这个技巧经常比求导快。
\end{explanationbox}
